% !TEX program = lualatex
\documentclass[a4paper,10pt,twocolumn]{article}
\usepackage{kaitoushu}
\renewcommand{\baselinestretch}{1.2}
\lhead{応用数学　練習問題（採点用）}
\problemrange{第3章 §1 練習問題1（p.\,90）}

\begin{document}
\thispagestyle{fancy}

\begin{center}
  {\Large \textbf{応用数学　第3章 §1　練習問題1（p.\,90）}}
\end{center}
\vspace{0.5em}

\noindent\textbf{\large【フーリエ展開】}
\vspace{0.3em}

\mondai{1}{
次の周期 $2$ の関数 $f(x)$ のフーリエ級数を求めよ．\\
$f(x) = \begin{cases} 1 & (-1 \leq x < 0) \\ 1-x & (0 \leq x < 1) \end{cases},\quad f(x+2) = f(x)$
}
\kotae{
$l=1$．\\[2pt]
$c_0 = \dfrac{1}{2}\!\left(\displaystyle\int_{-1}^{0}dx + \int_{0}^{1}(1-x)dx\right)
  = \dfrac{1}{2}\!\left(1+\dfrac{1}{2}\right) = \dfrac{3}{4}$．\\[4pt]
\textbf{$a_n$：}
$\displaystyle\int_{-1}^{0}\cos n\pi x\,dx
  = \!\left[\dfrac{\sin n\pi x}{n\pi}\right]_{-1}^{0}\! = 0$．\\[2pt]
$\displaystyle\int_{0}^{1}(1-x)\cos n\pi x\,dx$：$u=1-x,\;dv=\cos n\pi x\,dx$ より\\[2pt]
$= \!\left[\dfrac{(1-x)\sin n\pi x}{n\pi}\right]_0^1\!
  + \dfrac{1}{n\pi}\!\int_0^1\!\sin n\pi x\,dx$\\[2pt]
$= 0 + \dfrac{1}{n\pi}\!\left[-\dfrac{\cos n\pi x}{n\pi}\right]_0^1\!
  = \dfrac{1-(-1)^n}{n^2\pi^2}$．\\[4pt]
\textbf{$b_n$：}
$\displaystyle\int_{-1}^{0}\sin n\pi x\,dx
  = \!\left[-\dfrac{\cos n\pi x}{n\pi}\right]_{-1}^{0}\! = \dfrac{(-1)^n-1}{n\pi}$．\\[2pt]
$\displaystyle\int_{0}^{1}(1-x)\sin n\pi x\,dx$：$u=1-x,\;dv=\sin n\pi x\,dx$ より\\[2pt]
$= \!\left[-\dfrac{(1-x)\cos n\pi x}{n\pi}\right]_0^1\!
  - \dfrac{1}{n\pi}\!\int_0^1\!\cos n\pi x\,dx$\\[2pt]
$= \dfrac{1}{n\pi} - 0 = \dfrac{1}{n\pi}$．\\[2pt]
よって $b_n = \dfrac{(-1)^n-1}{n\pi}+\dfrac{1}{n\pi} = \dfrac{(-1)^n}{n\pi}$．
\[
  f(x) = \dfrac{3}{4} + \sum_{n=1}^\infty\!\left[\dfrac{1-(-1)^n}{n^2\pi^2}\cos n\pi x + \dfrac{(-1)^n}{n\pi}\sin n\pi x\right]
\]
}

\mondai{2}{
グラフが図で示されている次の関数のフーリエ級数を求めよ．\\
(1) 三角波：
$f(x) = \begin{cases} 2-2|x| & (-1 \leq x \leq 1) \\ 0 & (|x| > 1) \end{cases},\quad f(x+4) = f(x)$\\
(2) のこぎり波：
$f(x) = \begin{cases} -x-2 & (-2 \leq x < -1) \\ 0 & (-1 \leq x < 1) \\ 2-x & (1 \leq x < 2) \end{cases},\quad f(x+4) = f(x)$
}
\kotae{
(1) 周期 $4,\ l=2$．偶関数なので $b_n=0$．\\[2pt]
$c_0 = \dfrac{1}{4}\cdot 2\displaystyle\int_0^1(2-2x)dx
  = \dfrac{1}{2}\!\left[2x-x^2\right]_0^1\! = \dfrac{1}{2}$．\\[4pt]
$a_n = \displaystyle\int_0^1(2-2x)\cos\dfrac{n\pi x}{2}dx$：
$u=2-2x,\;dv=\cos\dfrac{n\pi x}{2}dx$ より\\[2pt]
$= \!\left[\dfrac{2(2-2x)\sin\frac{n\pi x}{2}}{n\pi}\right]_0^1\!
  + \dfrac{4}{n\pi}\!\int_0^1\!\sin\dfrac{n\pi x}{2}dx$\\[2pt]
$= 0 + \dfrac{4}{n\pi}\!\left[-\dfrac{2\cos\frac{n\pi x}{2}}{n\pi}\right]_0^1\!
  = \dfrac{8\!\left(1-\cos\dfrac{n\pi}{2}\right)}{n^2\pi^2}$．
\[
  f(x) = \dfrac{1}{2} + \sum_{n=1}^\infty \dfrac{8\!\left(1-\cos\dfrac{n\pi}{2}\right)}{n^2\pi^2}\cos\dfrac{n\pi x}{2}
\]
(2) 周期 $4,\ l=2$．奇関数なので $a_n=0$．\\[2pt]
$b_n = \displaystyle\int_1^2 (2-x)\sin\dfrac{n\pi x}{2}\,dx$：$u=2-x,\;dv=\sin\dfrac{n\pi x}{2}dx$ より\\[2pt]
$= \!\left[-\dfrac{2(2-x)\cos\frac{n\pi x}{2}}{n\pi}\right]_1^2\!
  - \dfrac{2}{n\pi}\!\int_1^2\!\cos\dfrac{n\pi x}{2}\,dx$\\[2pt]
$= \dfrac{2\cos\frac{n\pi}{2}}{n\pi}
  - \dfrac{2}{n\pi}\!\left[\dfrac{2\sin\frac{n\pi x}{2}}{n\pi}\right]_1^2\!
  = \dfrac{2\cos\frac{n\pi}{2}}{n\pi} + \dfrac{4\sin\frac{n\pi}{2}}{n^2\pi^2}$．
\[
  f(x) = \sum_{n=1}^\infty \!\left(\dfrac{2\cos\frac{n\pi}{2}}{n\pi} + \dfrac{4\sin\frac{n\pi}{2}}{n^2\pi^2}\right)\sin\dfrac{n\pi x}{2}
\]
}

\mondai{3}{
次の関数について以下の問いに答えよ．
$f(x) = x^2\ (-1 \leq x < 1),\ f(x+2) = f(x)$\\
(1) フーリエ級数を求めよ．\\
(2) $x=0$ および $x=1$ とおくことにより
$1 - \dfrac{1}{2^2} + \dfrac{1}{3^2} - \cdots = \dfrac{\pi^2}{12}$，
$1 + \dfrac{1}{2^2} + \dfrac{1}{3^2} + \cdots = \dfrac{\pi^2}{6}$ を導け．
}
\kotae{
(1) 偶関数なので $b_n=0$．\\[2pt]
$c_0 = \dfrac{1}{2}\displaystyle\int_{-1}^{1}x^2dx = \dfrac{1}{3}$．\\[4pt]
$a_n = 2\displaystyle\int_0^1 x^2\cos n\pi x\,dx$（部分積分2回）：\\[2pt]
$u=x^2,\;dv=\cos n\pi x\,dx$ より\\[2pt]
$= 2\!\left\{\!\left[\dfrac{x^2\sin n\pi x}{n\pi}\right]_0^1\!
  - \dfrac{2}{n\pi}\!\int_0^1\!x\sin n\pi x\,dx\right\}$\\[2pt]
$= \dfrac{-4}{n\pi}\displaystyle\int_0^1\!x\sin n\pi x\,dx$．\\[4pt]
$u=x,\;dv=\sin n\pi x\,dx$ より\\[2pt]
$\displaystyle\int_0^1\!x\sin n\pi x\,dx
  = \!\left[-\dfrac{x\cos n\pi x}{n\pi}\right]_0^1\!
  + \dfrac{1}{n\pi}\!\int_0^1\!\cos n\pi x\,dx
  = \dfrac{(-1)^{n+1}}{n\pi}$．\\[4pt]
よって $a_n = \dfrac{-4}{n\pi}\cdot\dfrac{(-1)^{n+1}}{n\pi} = \dfrac{4(-1)^n}{n^2\pi^2}$．
\[
  f(x) = \dfrac{1}{3} + \sum_{n=1}^\infty \dfrac{4(-1)^n}{n^2\pi^2}\cos n\pi x
\]
(2) $x=0$ で $f$ は連続で $f(0)=0$ より\\[2pt]
$0 = \dfrac{1}{3}+\dfrac{4}{\pi^2}\displaystyle\sum_{n=1}^\infty\dfrac{(-1)^n}{n^2}$，すなわち
$\displaystyle\sum_{n=1}^\infty\dfrac{(-1)^{n+1}}{n^2} = \dfrac{\pi^2}{12}$．\\[4pt]
$x=1$ で $f$ は連続で $f(1)=1$ より\\[2pt]
$\cos n\pi=(-1)^n$ より $1 = \dfrac{1}{3}+\dfrac{4}{\pi^2}\displaystyle\sum_{n=1}^\infty\dfrac{1}{n^2}$，\\[2pt]
すなわち $\displaystyle\sum_{n=1}^\infty\dfrac{1}{n^2} = \dfrac{\pi^2}{6}$．
}

\mondai{4}{
次の関数 $f(x)$ の複素フーリエ級数を求めよ．\\
$f(x) = e^x\ (-\pi/2 \leq x < \pi/2),\ f(x+\pi) = f(x)$
}
\kotae{
周期 $\pi,\ l=\pi/2$．\\[4pt]
\textbf{$n=0$：}
$c_0 = \dfrac{1}{\pi}\displaystyle\int_{-\pi/2}^{\pi/2}e^x\,dx
  = \dfrac{1}{\pi}\bigl[e^x\bigr]_{-\pi/2}^{\pi/2}
  = \dfrac{e^{\pi/2}-e^{-\pi/2}}{\pi} = \dfrac{2\sinh(\pi/2)}{\pi}$．\\[4pt]
\textbf{$n\neq 0$：}
\[
  c_n = \dfrac{1}{\pi}\int_{-\pi/2}^{\pi/2}\!e^{(1-2in)x}\,dx
      = \dfrac{1}{\pi(1-2in)}\bigl[e^{(1-2in)x}\bigr]_{-\pi/2}^{\pi/2}
\]
$e^{(1-2in)\pi/2}=(-1)^n e^{\pi/2}$，$e^{-(1-2in)\pi/2}=(-1)^n e^{-\pi/2}$ より\\[2pt]
$\bigl[e^{(1-2in)x}\bigr]_{-\pi/2}^{\pi/2} = 2(-1)^n\sinh\dfrac{\pi}{2}$，\\[4pt]
$c_n = \dfrac{2(-1)^n\sinh(\pi/2)}{\pi(1-2in)}$．
分母有理化 $\left(\cdot\dfrac{1+2in}{1+2in}\right)$：
\[
  c_n = \dfrac{2(-1)^n(1+2in)\sinh(\pi/2)}{\pi(1+4n^2)}
\]
\[
  f(x) = c_0 + \sum_{\substack{n=-\infty \\ n\neq 0}}^{\infty} c_n\,e^{2inx}
\]
}

\end{document}
